\begin{ZhChapter}

\chapter{實驗與結果分析}

本章旨在評估 CBR-APIKGAgent 於多步驟 API 任務中的任務完成能力與錯誤修復表現。透過 AppWorld 任務環境進行實驗，觀察完整系統之整體表現，並進一步分析原子化經驗重用與動態局部修復機制各自對任務完成率造成的影響。

由於本研究的核心假設為「由訓練資料成功任務軌跡萃取之原子化經驗，可作為相似任務結構或錯誤情境下的策略參考」，以及「局部修復可在部分失敗情境下協助代理人由執行錯誤中恢復」，因此本章將分別從整體效能、消融實驗與案例分析等面向進行分析。實驗結果以實際執行紀錄與 AppWorld 評估結果為依據，並搭配可取得完整軌跡的案例，討論兩項機制如何在不同階段共同影響任務完成，以及完整系統目前呈現的限制。

\section{研究問題}

本研究設計四項研究問題（Research Questions, RQs），分別以 RQ1 至 RQ4 表示，以對應第一章所提出之研究目的：

\begin{enumerate}
    \item \textbf{RQ1：CBR-APIKGAgent 相較於既有方法在 AppWorld 任務中的整體表現為何？}

    此問題用於評估完整系統在 AppWorld \texttt{test\_normal} 上的整體任務完成能力。

    \item \textbf{RQ2：原子化經驗重用機制對任務完成率造成何種影響？}

    此問題用於透過消融實驗分析原子化經驗重用機制對任務完成結果的影響。

    \item \textbf{RQ3：動態局部修復機制對任務完成率造成何種影響？}

    此問題用於透過消融實驗分析動態局部修復機制對任務完成結果的影響。

    \item \textbf{RQ4：原子化經驗與動態局部修復如何共同影響多步驟 API 任務的完成？}

    此問題延續消融實驗結果，透過可取得完整軌跡的成功與限制案例，觀察原子化經驗與動態局部修復分別在規劃與執行階段的作用，以及兩者如何共同影響任務完成。
\end{enumerate}

\section{實驗設置}

\subsection{資料集與任務範圍}

本研究於 AppWorld \cite{trivedi2024appworld} 環境中進行實驗。AppWorld 提供多應用程式 API 任務環境，任務完成與否由最終環境狀態是否符合預期條件判定。因此，代理人不僅需要產生合理文字回覆，也必須透過 API 呼叫與資料處理正確改變或查詢環境狀態。

正式實驗以 AppWorld \texttt{test\_normal} 作為主要評估資料集，共包含 168 個任務，並以本研究提出之 CBR-APIKGAgent 作為評估系統。各項實驗所使用的任務範圍將於對應實驗小節中說明。

在每一項實驗設定中，各測試任務僅採用一次正式執行結果；局部修復、重新規劃與執行重試皆發生於該次任務執行期間，並受後述控制設定限制。任務結束後，AppWorld evaluation tests 的結果僅用於計算評估指標，不會回饋給代理人、寫入正式經驗庫，或作為調整後重新執行同一測試任務的依據。

\subsection{評估指標}

依據 AppWorld 官方評估方式，AppWorld 採用程式化且以最終狀態為基礎的 evaluation tests 判斷任務是否完成，並以任務目標完成率（Task Goal Completion, TGC）與情境目標完成率（Scenario Goal Completion, SGC）作為主要指標 \cite{trivedi2024appworld}。

\begin{enumerate}
    \item \textbf{任務目標完成率（TGC）：}TGC 表示通過所有 AppWorld evaluation tests 的任務比例。對於單一任務而言，只有當該任務的所有評估測試皆通過時，該任務才計為成功。其計算方式如式 \ref{eq:tgc} 所示：

    \begin{equation}
        \mathrm{TGC} = \frac{N_{\mathrm{success}}}{N_{\mathrm{total}}} \times 100\%
        \label{eq:tgc}
    \end{equation}

    其中，\(N_{\mathrm{success}}\) 代表通過所有 evaluation tests 的任務數量，\(N_{\mathrm{total}}\) 代表評估任務總數。

    \item \textbf{情境目標完成率（SGC）：}
    SGC 以 AppWorld 的 task scenario 為單位計算。只有同一 scenario 產生之所有任務皆通過 evaluation tests 時，該 scenario 才計為成功。此指標用於衡量代理人是否能在同一任務情境的不同需求與初始狀態變化下穩定完成目標。其計算方式如式 \ref{eq:sgc} 所示：

    \begin{equation}
        \mathrm{SGC} = \frac{N_{\mathrm{success\ scenarios}}}{N_{\mathrm{total\ scenarios}}} \times 100\%
        \label{eq:sgc}
    \end{equation}
\end{enumerate}

\subsection{實驗環境與主要控制設定}

本研究實驗之硬體與軟體環境如表 \ref{tab:hardware_software_environment} 所示。由於本研究主要透過雲端大型語言模型服務進行推論，本地端硬體主要負責執行 AppWorld 環境、代理人控制流程、API 模擬環境與實驗紀錄整理。

\begin{table}[htbp]
    \centering
    \caption{硬體與軟體環境}
    \label{tab:hardware_software_environment}
    \small
    \renewcommand{\arraystretch}{1.15}
    \begin{tabularx}{0.82\textwidth}{
        >{\raggedright\arraybackslash}p{0.3\textwidth}
        >{\raggedright\arraybackslash}X
    }
        \hline
        \textbf{項目} & \textbf{設定} \\
        \hline
        作業系統 & Ubuntu 22.04.5 LTS（WSL2） \\
        處理器 & 12th Gen Intel(R) Core(TM) i7-12700 \\
        記憶體 & 16 GB \\
        Python 版本 & Python 3.12 \\
        AppWorld 版本 & 0.1.3 \\
        \hline
    \end{tabularx}
\end{table}

除上述執行環境外，實驗亦需固定語言模型、推論溫度、執行預算、經驗檢索設定與局部修復限制等主要控制項。本研究採用 Gemini 2.5 Flash 作為大型語言模型 \cite{geminiteam2025gemini25}；表 \ref{tab:experiment_environment} 列出本研究實驗中需要固定或比較的重點設定。

\begin{table}[htbp]
    \centering
    \caption{實驗主要控制設定}
    \label{tab:experiment_environment}
    \small
    \renewcommand{\arraystretch}{1.15}
    \begin{tabularx}{0.82\textwidth}{
        >{\raggedright\arraybackslash}p{0.3\textwidth}
        >{\raggedright\arraybackslash}X
    }
        \hline
        \textbf{項目} & \textbf{設定} \\
        \hline
        大型語言模型 & \texttt{gemini-2.5-flash} \\
        溫度參數 & 判斷 0；規劃、重新規劃與執行 0.2；API 評估 0.1 \\
        任務執行時間限制 & 3600 秒／題 \\
        執行流程內重新嘗試上限 & 3 次／題 \\
        局部修復上限 & 3 次／原始失敗步驟 \\
        修復鏈深度上限 & 2 步驟／原始失敗步驟 \\
        經驗檢索相似度門檻 & 0.35 \\
        初始規劃經驗檢索 & 規劃導向至多 3 筆 \\
        重新規劃經驗檢索 & 規劃導向至多 1 筆；修復導向至多 1 筆 \\
        局部修復經驗檢索 & 修復導向至多 2 筆 \\
        \hline
    \end{tabularx}
\end{table}

\section{實驗一：整體效能比較}
\label{sec:overall_performance}

本節對應 RQ1，目的在於呈現 CBR-APIKGAgent 在 AppWorld \texttt{test\_normal} 上的整體任務完成率，並與既有方法進行比較。本實驗使用完整 \texttt{test\_normal} 168 個任務作為評估範圍，並依 AppWorld 官方評估結果計算 TGC 與 SGC。若任務執行後未產生完整評估報告或出現異常執行情形，該任務在統計中計為失敗。

為呈現本研究方法於 AppWorld 基準中的相對位置，表 \ref{tab:overall_performance} 列出既有方法於 \texttt{test\_normal} 上的結果。表中 GPT-4o 與 LLaMA3-70B 設定下之 ReAct、PlanExec、FullCodeRefl 與 IPFunCall 結果取自 AppWorld 原論文 \cite{trivedi2024appworld}；Qwen2.5-32B 設定下之 ReAct，以及 Gemini 2.5 Flash 設定下之 ReAct、FullCodeRefl 與 APIKGAgent 結果取自 APIKGAgent 原研究整理 \cite{wu2025apikgagent}；IBM CUGA 之結果取自其原研究報告 \cite{shlomov2026benchmarks}；LOOP 之結果取自本研究整理時可取得之 AppWorld 官方 leaderboard \cite{appworld_leaderboard}，其方法背景則參考 LOOP 原研究 \cite{chen2025loop}。其中，leaderboard 數據以 2026 年 7 月 13 日存取結果為準。CBR-APIKGAgent 為本研究於相同 \texttt{test\_normal} 任務集合下之實驗結果。

由於不同方法可能採用不同模型、示範資料、訓練策略或執行預算，因此表 \ref{tab:overall_performance} 主要作為 benchmark 相對表現參考，而非嚴格控制變因下的直接比較。各機制對本研究系統造成的影響，將於後續消融實驗中以相同模型與相同任務集合進一步分析。

\begin{table}[htbp]
    \centering
    \caption{各方法於 AppWorld \texttt{test\_normal} 之整體效能比較}
    \label{tab:overall_performance}
    \small
    \renewcommand{\arraystretch}{1.15}
    \begin{tabularx}{0.8\textwidth}{
        >{\raggedright\arraybackslash}X
        >{\raggedright\arraybackslash}p{0.22\textwidth}
        >{\centering\arraybackslash}p{0.12\textwidth}
        >{\centering\arraybackslash}p{0.12\textwidth}
    }
        \hline
        \textbf{Method} & \textbf{Model} & \textbf{TGC (\%)} & \textbf{SGC (\%)} \\
        \hline
        ReAct & \multirow{4}{*}{GPT-4o} & 48.8 & 32.1 \\
        PlanExec &  & 44.6 & 23.2 \\
        FullCodeRefl &  & 33.9 & 26.8 \\
        IPFunCall &  & 32.1 & 16.1 \\
        \hline
        IBM CUGA & GPT-4.1 & 73.2 & 62.5 \\
        \hline
        ReAct & \multirow{3}{*}{LLaMA3-70B} & 20.8 & 8.9 \\
        PlanExec &  & 8.9 & 1.8 \\
        FullCodeRefl &  & 24.4 & 17.9 \\
        \hline
        ReAct & \multirow{2}{*}{Qwen2.5-32B} & 39.2 & 18.6 \\
        LOOP &  & 72.6 & 53.6 \\
        \hline
        ReAct & \multirow{4}{*}{Gemini 2.5 Flash} & 50.6 & 28.6 \\
        FullCodeRefl &  & 36.9 & 26.8 \\
        APIKGAgent &  & 41.7 & 17.9 \\
        \textbf{CBR-APIKGAgent} &  & \textbf{51.8} & \textbf{35.7} \\
        \hline
    \end{tabularx}
\end{table}

由表 \ref{tab:overall_performance} 可知，CBR-APIKGAgent 於 \texttt{test\_normal} 上達成 51.8\% 的 TGC 與 35.7\% 的 SGC。在表中採用 Gemini 2.5 Flash 的方法裡，本研究方法之兩項指標皆為最高，並高於同模型的 ReAct、FullCodeRefl 與 APIKGAgent；但與 LOOP 及 IBM CUGA 相比仍有差距。

從方法設計來看，LOOP 透過強化學習直接在 AppWorld 環境中調整代理人的行為，使 API 文件查閱、行動選擇與錯誤恢復策略反映於模型的互動行為；本研究則不調整模型權重，而是在推論階段檢索原子化經驗，並將其加入規劃與修復脈絡以引導模型重用既有策略。因此，本研究方法的效果仍會受到經驗檢索結果、任務限制契合度及模型對經驗內容之理解所影響，此項學習方式的差異可能是 LOOP 取得較高完成率的原因之一。另一方面，CUGA 採用階層式規劃器與執行器架構，由持久化任務紀錄維持步驟、變數綁定與重新規劃狀態，再將子任務交由專用代理人處理；本研究則著重於在 APIKGAgent 流程中加入經驗重用與失敗後的局部修復。從本研究的限制案例可觀察到，當任務涉及多階段資料篩選或搜尋候選結果時，原始任務條件與後續操作目標之間的對應仍可能產生偏差，因此，跨步驟狀態維持方式的差異可能是 CUGA 取得較高完成率的原因之一。

然而，LOOP、CUGA 與本研究使用的基礎模型、學習方式及執行架構並不相同，本研究亦未在相同模型與執行預算下重現上述方法。因此，前述差異應視為根據各方法設計與本研究案例觀察提出的可能解釋，尚不能直接判定為效能差距的單一原因。針對 RQ1，本實驗顯示 CBR-APIKGAgent 在相同模型條件下改善了整體任務完成表現，但尚未達到 AppWorld 公開結果中的最佳水準。至於兩項核心機制各自的貢獻，將於下一節透過消融實驗進一步分析。

\section{實驗二：消融實驗與案例分析}

本節對應 RQ2、RQ3 與 RQ4。首先，透過消融實驗分別分析原子化經驗重用機制與動態局部修復機制對系統表現造成的影響，以回答 RQ2 與 RQ3。其次，本節延續消融實驗結果，透過可取得完整軌跡的成功與限制案例，觀察兩項機制如何在規劃與執行階段共同影響任務完成，以回答 RQ4。

\subsection{實驗組別}

消融實驗共包含四種設定，用以分別觀察原子化經驗與局部修復對任務完成結果的影響。四種設定皆於 AppWorld \texttt{test\_normal} 的 168 個任務進行評估，並採相同計分方式。除上述機制差異外，其餘模型、時間限制、重試上限與評估方式皆保持一致。

\subsection{消融實驗結果}

正式消融實驗以基礎系統作為比較基準，分別加入原子化經驗、動態局部修復，以及同時啟用兩項機制，以觀察各項機制對任務完成率的影響。四種設定皆使用相同任務集合，並以 TGC 與 SGC 作為比較指標，其中表中的 $\Delta$ 代表相較於 Baseline System 的提升幅度，單位為 pp。實驗結果顯示，僅啟用原子化經驗或僅啟用局部修復時，兩項指標皆高於基礎系統；同時啟用兩項機制的完整系統則取得四種設定中最高的 TGC 與 SGC。各組實際結果整理如表 \ref{tab:ablation_results} 所示。

\begin{table}[htbp]
    \centering
    \caption{消融實驗結果}
    \label{tab:ablation_results}
    \small
    \renewcommand{\arraystretch}{1.15}
    \setlength{\tabcolsep}{4pt}
    \begin{tabularx}{1\textwidth}{
        >{\raggedright\arraybackslash}X
        >{\centering\arraybackslash}p{0.13\textwidth}
        >{\centering\arraybackslash}p{0.13\textwidth}
        >{\centering\arraybackslash}p{0.13\textwidth}
        >{\centering\arraybackslash}p{0.13\textwidth}
    }
        \hline
        \textbf{Configuration} & \textbf{TGC (\%)} & \textbf{\ensuremath{\Delta}TGC (pp)} & \textbf{SGC (\%)} & \textbf{\ensuremath{\Delta}SGC (pp)} \\
        \hline
        Baseline System & 37.5 & --- & 14.3 & --- \\
        Atomic Experience Only & 43.5 & +6.0 & 25.0 & +10.7 \\
        Local Repair Only & 49.4 & +11.9 & 28.6 & +14.3 \\
        Full System & \textbf{51.8} & \textbf{+14.3} & \textbf{35.7} & \textbf{+21.4} \\
        \hline
    \end{tabularx}
\end{table}

為補充整體 TGC 與 SGC 指標，本研究進一步比較不同消融設定下的 task-level 評估結果差異。表 \ref{tab:task_level_ablation_comparison} 列出三類與研究問題相關的正向差集，用以觀察各機制可能帶來改善的任務集合。

\begin{table}[htbp]
    \centering
    \caption{消融設定之任務層級正向差集比較}
    \label{tab:task_level_ablation_comparison}
    \small
    \renewcommand{\arraystretch}{1.15}
    \begin{tabularx}{0.9\textwidth}{
        >{\raggedright\arraybackslash}X
        >{\centering\arraybackslash}p{0.22\textwidth}
    }
        \hline
        \textbf{Task-level outcome comparison} & \textbf{Number of tasks} \\
        \hline
        Baseline failed, AE-only succeeded & 19 \\
        Baseline failed, Local Repair-only succeeded & 32 \\
        Both single-mechanism settings failed, Full System succeeded & 8 \\
        \hline
    \end{tabularx}
\end{table}

各比較條件分別統計，前兩類可能包含相同任務。表 \ref{tab:task_level_ablation_comparison} 僅呈現不同消融設定下的任務評估結果差異，用以觀察各機制可能帶來改善的任務集合；兩項機制的實際作用方式仍需配合執行軌跡分析。因此，本表不應被解讀為完整的任務轉移矩陣，也不直接將任務成功歸因於單一機制。

後續將先分析原子化經驗與局部修復各自的影響，再透過完整系統與可取得完整軌跡的案例，討論兩項機制如何在規劃與執行階段共同影響任務完成。

\subsection{單一機制效益分析}

本節分別比較僅原子化經驗與僅局部修復兩種設定相較於基礎系統的差異，以回答 RQ2 與 RQ3。

\subsection*{原子化經驗之影響}

由表 \ref{tab:ablation_results} 可知，基礎系統之 TGC 為 37.5\%，SGC 為 14.3\%。僅原子化經驗設定之 TGC 為 43.5\%，SGC 為 25.0\%，相較於基礎系統分別提升 6.0 與 10.7 個百分點。此結果表示，在未啟用動態局部修復的情況下，原子化經驗設定仍呈現較高的任務層級與情境層級完成率。

進一步從任務層級比較來看，表 \ref{tab:task_level_ablation_comparison} 顯示「基礎系統失敗、僅原子化經驗成功」的任務共有 19 題。此差集表示，在部分任務中，加入原子化經驗後可觀察到原本未通過評估的任務轉為通過；然而，這些任務的改善情形仍需配合執行軌跡檢視，不能僅由差集結果直接歸因於單一機制。

以下以一個 Venmo payment request 修正任務作為代表性案例，說明其中一種可觀察到的改善情形。此案例用於呈現原子化經驗可能如何協助系統在規劃階段建立目標集合與篩選流程，並不代表所有僅原子化經驗成功的任務皆具有相同原因。

\noindent\textbf{範例：}
\begin{itemize}
    \item \textbf{案例任務：}使用者需修正前一晚向 roommates 發出的 Venmo payment requests，刪除原本金額少 10 美元的請求，並以相同描述與收款對象建立修正後的新請求。
    \item \textbf{消融差異：}此任務在基礎系統中未通過評估，但在僅原子化經驗設定中成功。
    \item \textbf{基礎系統觀察：}基礎系統將 roommate 身分視為缺失資訊，反覆要求使用者提供 roommates 的 names 或 email addresses，因而未能從既有資料中建立可用的目標集合。
    \item \textbf{機制作用：}僅原子化經驗設定的軌跡中，系統在初始規劃與後續重新規劃過程中，檢索到與寫入前目標集合建立、既有資料比對、目標篩選，以及使用 stable identifiers 維持操作對象一致性相關的策略經驗。此類經驗可能與後續規劃中的目標集合建立行為相關；在該軌跡中，系統將 roommates 這類隱含對象條件轉換為可查詢的 contact relationship 條件，先透過 phone contacts 查詢並以 roommate 關係取得 roommates 的 email；接著登入 Venmo、取得 sent payment requests，依 roommates email 與前一晚時間條件篩選待修正 requests，最後刪除原請求並建立金額增加 10 美元的新請求。
    \item \textbf{影響：}此代表性軌跡顯示，原子化經驗可能有助於在規劃階段建立目標集合與篩選流程，特別是將任務中的隱含對象條件轉換為可執行的資料取得步驟。不過，此案例不代表所有僅原子化經驗改善的任務皆具有相同原因。
\end{itemize}
\subsection*{動態局部修復之影響}

僅局部修復設定之 TGC 為 49.4\%，SGC 為 28.6\%，相較於基礎系統分別提升 11.9 與 14.3 個百分點。此結果表示，在不加入原子化經驗的情況下，執行失敗後的局部計畫修正仍可能使部分原本失敗的任務恢復並完成。進一步從任務層級比較來看，表 \ref{tab:task_level_ablation_comparison} 顯示「基礎系統失敗、僅局部修復成功」的任務共有 32 題。此差集可用於觀察局部修復可能帶來改善的任務集合，但各任務的實際改善方式仍需搭配執行軌跡分析，不能直接概括為同一種失敗原因。

從已檢視的代表性軌跡來看，局部修復可觀察到對不同類型執行問題的補救行為，例如授權缺失、必要前置資料不足、API 選擇不適合或查詢條件不精確等。這些例子顯示，動態局部修復可能與執行過程中利用局部錯誤證據修正當前失敗點的行為相關，使系統有機會在不重新規劃整個任務的情況下恢復執行。

\noindent\textbf{範例：}
\begin{itemize}
    \item \textbf{案例任務：}使用者需更新 Spotify 播放清單，依據兄弟姊妹在手機訊息中的建議調整歌曲內容。
    \item \textbf{消融差異：}此任務在基礎系統中失敗，但在僅局部修復設定中成功。
    \item \textbf{基礎系統觀察：}基礎系統能確認 sibling 關係，查到三位 sibling contacts，並取得包含歌曲新增與移除建議的 text messages。然而，後續流程步驟在辨識目標 Spotify playlist 時，未能將欲辨識的 roadtrip playlist 對應到實際的 Road Trip 播放清單；由於該步驟在未找到相符播放清單後即進入停止流程，後續歌曲搜尋、新增與移除操作皆未執行，因此 evaluation tests 中沒有產生預期的 playlist song 變更。
    \item \textbf{局部修復軌跡：}僅局部修復設定的初始流程曾使用 family 關係查詢 contacts 以尋找 family members 或 siblings，但結果為空，導致後續從空結果中擷取 sibling phone numbers 的步驟無法繼續。系統因此觸發局部修復，新增 phone contacts 查詢作為修復步驟，並改用 sibling 關係查詢。
    \item \textbf{機制作用：}在該軌跡中，修復步驟取得三位 sibling contacts 後，系統繼續搜尋 text messages、解析歌曲新增與移除建議、取得 playlist、搜尋歌曲，最後透過 Spotify playlist API 執行歌曲新增與移除，並通過 evaluation tests。
    \item \textbf{影響：}此代表性案例顯示，動態局部修復可能有助於在執行流程中根據錯誤回饋補足缺失的前置資料，協助任務從局部失敗點恢復執行。不過，此案例中的局部修復點與基礎系統的失敗點並不相同，因此只能說明一種可觀察到的修復行為，不能推論所有僅局部修復成功的任務皆具有相同原因。
\end{itemize}

整體而言，原子化經驗與動態局部修復在單獨啟用時皆高於基礎系統。針對 RQ2，僅原子化經驗設定使 TGC 與 SGC 分別提升 6.0 與 10.7 個百分點，代表性案例則顯示其可能透過提供資料取得、條件篩選與操作順序等策略參考，改善部分任務的規劃完整性。針對 RQ3，僅局部修復設定使 TGC 與 SGC 分別提升 11.9 與 14.3 個百分點，代表性案例則顯示其可能利用執行階段的局部錯誤證據修正失敗步驟，提升部分任務的恢復能力。在本組實驗中，僅局部修復設定之兩項指標皆高於僅原子化經驗設定；代表性案例顯示，兩項機制可分別在規劃與執行恢復階段產生作用，並對任務完成結果產生不同層次的影響。

\subsection{完整系統與代表性案例分析}
\label{sec:ablation_case_analysis}

本節對應 RQ4，目的在於說明原子化經驗與動態局部修復可能如何共同影響任務完成。相較於前一節以消融數據分析兩項機制各自對任務完成率的影響，本節先透過可取得完整軌跡的成功案例，觀察兩者在規劃與執行階段可能形成的互補關係；再透過限制案例，討論目前完整系統仍可能面臨的限制。

完整系統同時啟用原子化經驗與動態局部修復，其 TGC 為 51.8\%，SGC 為 35.7\%，為四種消融設定中最高。相較於基礎系統，完整系統之 TGC 提升 14.3 個百分點，SGC 提升 21.4 個百分點；相較於僅原子化經驗設定，完整系統之 TGC 與 SGC 分別提升 8.3 與 10.7 個百分點；相較於僅局部修復設定，完整系統之 TGC 與 SGC 分別提升 2.4 與 7.1 個百分點。

從任務層級差集來看，表 \ref{tab:task_level_ablation_comparison} 顯示「僅原子化經驗設定與僅局部修復設定皆失敗，但完整系統成功」的任務共有 8 題。此結果表示，部分任務僅在兩項機制同時啟用時通過評估，可作為觀察兩項機制可能互補的起點；但此差集本身不能直接證明兩項機制在所有任務中皆具有協同作用，仍需搭配具體執行軌跡分析。

\subsection*{互補案例分析}

根據可取得完整軌跡的完整系統成功案例可觀察到，原子化經驗可能在規劃階段提供目標集合建立與篩選策略的參考；若執行過程仍出現授權或前置資訊缺失，動態局部修復則可能依據當下錯誤補足必要步驟。此類案例顯示，兩項機制可能與不同階段的行為變化相關，並在部分任務中共同影響任務完成。

\noindent\textbf{範例：}
\begin{itemize}
    \item \textbf{案例任務：}使用者需將手機聯絡人中的 coworkers 與 friends 加入 Venmo friends，但僅限於目前尚未是 Venmo friends 的對象。
    \item \textbf{消融差異：}此任務在僅原子化經驗設定與僅局部修復設定中皆未通過評估，但在完整系統中成功。
    \item \textbf{單一機制觀察：}僅原子化經驗設定雖執行至新增好友步驟，但在將手機聯絡人對應到 Venmo 使用者時，搜尋結果包含了不屬於原始 coworkers 或 friends 的額外對象；後續篩選又未充分扣回「手機聯絡人中的 coworkers 或 friends，且尚未是 Venmo friends」這一任務條件，因此最後新增 13 位使用者，新增集合超出評估預期。僅局部修復設定則可處理 Venmo 授權錯誤與部分漏處理的查詢結果，但修復後仍沿用已被擴大的候選集合，因此最後新增 5 位使用者，仍未符合評估條件。
    \item \textbf{完整系統軌跡：}完整系統先從 phone contacts 取得候選對象，篩選 relationships 為 coworker 或 friend 的 contacts，形成候選好友 email 集合。接著系統嘗試透過 Venmo friends 查詢取得既有 Venmo friends，但遇到 401 authentication error；動態局部修復觸發後，新增 Venmo login 作為修復步驟，取得 Venmo access token 後回到原本的 friends 查詢流程。
    \item \textbf{機制作用：}在該軌跡中，完整系統於修復後抽取既有 Venmo friend emails，再與候選好友 email 集合做差集，形成待新增 email 集合，最後只新增應加入的 7 位使用者。
    \item \textbf{影響：}此代表性案例顯示，原子化經驗相關流程可能有助於規劃階段的目標集合建立與既有狀態比對，動態局部修復則可能在執行階段處理 Venmo 授權錯誤，協助原流程繼續執行。不過，此案例僅說明部分完整系統成功任務中可觀察到的互補關係，不能概括為所有任務皆具有相同原因。
\end{itemize}
\subsection*{完整系統的限制分析}

雖然完整系統在整體 TGC 與 SGC 上皆為四種設定中最高，但在個別任務中，並不保證一定優於僅啟用單一機制的設定。這類情況不宜直接解讀為原子化經驗與動態局部修復彼此干擾；較合理的解釋是，完整系統的效果仍會受到任務結構、搜尋式 API 語意，以及後續目標集合維持方式影響。

以下案例與前述成功案例同樣涉及 Venmo friends 與 phone contacts 的對應，但呈現的是另一種情況：完整系統雖能處理部分執行錯誤，仍可能在目標集合維持上遺漏應新增對象。此例僅呈現其中一種失敗型態，並非所有完整系統失敗案例的共同原因。

\noindent\textbf{範例：}
\begin{itemize}
    \item \textbf{案例任務：}使用者需將 Venmo 上的朋友名單調整為與手機聯絡人中的朋友一致，必要時執行新增與移除操作。
    \item \textbf{消融差異：}此任務中，基礎系統與僅局部修復設定皆成功，但完整系統失敗。
    \item \textbf{完整系統觀察：}完整系統雖在初始計畫中納入手機聯絡人中 friend 關係的篩選，但後續使用 venmo.search\_users 將 phone friends 對應到 Venmo users 時，未能完整保留所有應新增對象。結果中實際新增的 Venmo friends 少於評估預期，表示應新增對象未被完整納入操作目標集合。
    \item \textbf{局部修復限制：}局部修復較可能處理由執行結果或語意判斷所辨識的步驟失敗，例如登入、授權或缺少查詢步驟等問題；但當問題發生在上游目標集合建立或候選對象維持時，後續修復即使被觸發，也未必能完整恢復原始任務條件。
    \item \textbf{影響：}此例顯示，完整系統雖能結合經驗檢索與局部修復，但在涉及搜尋式 API 與寫入操作的任務中，仍需要明確維持「原始任務條件 → 候選結果 → 實際操作目標」之間的對應關係。若應新增對象在候選集合建立過程中被遺漏，後續局部修復即使能補足授權或查詢步驟，也未必能使操作目標集合恢復完整。此限制案例不代表完整系統失敗皆源於相同原因。
\end{itemize}
從系統流程來看，API 呼叫完成後並非僅依據是否發生執行錯誤判斷步驟結果，系統亦會利用任務描述、步驟目標、API 參數、回傳結果與可取得的相依資料進行語意判斷。然而，局部修復與重新規劃仍以目前步驟先被判定為失敗為前提。若搜尋 API 回傳的候選集合未完整保留原始任務條件，或原始任務限制未完整呈現於後續步驟及相依資訊中，該步驟仍可能被判定為成功。此時，若後續流程亦未產生足以暴露偏差的執行回饋，不完整或不一致的操作目標可能持續傳遞，而不一定能透過後續局部修復或重新規劃完全修正。

因此，針對 RQ4，消融結果與代表性案例可觀察到，原子化經驗多作用於規劃階段的任務結構與資料流策略，動態局部修復則多作用於執行失敗後的局部計畫調整；兩者因作用階段不同，可能在部分任務中形成互補，並與完整系統取得四種設定中的最佳整體結果相互呼應。然而，兩項機制的結合並非在所有任務中皆能帶來正向效果；當檢索經驗未能維持任務限制，或上游目標集合已產生偏差時，局部修復未必能恢復正確的任務條件。此結果提示未來可從三個方向持續改善：在經驗使用時納入任務限制、細化局部修復與重新規劃的適用範圍，以及建立更可靠的經驗累積機制，以維持任務條件、候選結果與實際操作目標之間的對應關係。

\section{實驗效度限制}

為降低實驗結果受到非研究變因影響，本研究於消融實驗中使用相同的任務集合、基礎模型、評估指標與主要執行設定，並僅改變原子化經驗與動態局部修復兩項機制是否啟用，以比較各機制對任務完成率的影響。然而，大型語言模型仍具有生成不確定性，即使在相同設定下，單次執行結果仍可能受到模型輸出差異、API 呼叫順序或中間推理結果影響。因此，本研究結果可反映各設定在相同控制條件下的相對表現；未來若進一步進行多次重複執行，則可更細緻估計模型輸出變動對結果的影響。

此外，本研究主要於 AppWorld \texttt{test\_normal} 上進行評估。AppWorld 提供標準化的多應用程式 API 環境與程式化 evaluation tests，使不同方法能在相同任務定義下比較任務完成情形。然而，真實 API 環境可能包含更動態的資料狀態、權限限制、文件不一致、非同步操作與外部服務錯誤等因素。因此，本研究結果能說明 CBR-APIKGAgent 在標準化 API 任務基準中的表現與限制，但其效果是否能完全推廣至真實 API 服務或其他任務環境，仍需進一步驗證。

\end{ZhChapter}
